
\section{Весовая эвристика}

Кажется достаточно разумным, чтобы уменьшить итоговую высоту поднятия, нужно подвешивать множество меньшего размера к большему.
Такая эвристика называется \textit{весовой}. Нам необходимо поддерживать размер текущего множества, что достаточно тривиально:
\begin{enumerate}
    \item При \texttt{MAKE\_SET} будет создавать новый узел и в счётчик ставить 1.
    \item При \texttt{UNITE} будем просто прибавлять к счётчику размера большего дерева счётчик меньшего.
\end{enumerate}

\Def Пусть $x$ -- элемент некоторого множества в СНМ. Тогда $s(x)$, по определению, положим размером данного множества, в котором лежит $x$.

\Def Пусть $x$ и $y$ -- элементы некоторого множества в СНМ. Тогда $x \cup y$, по определению, результат объединения множеств, в которых они лежат. $s(x \cup y)$ -- размер результрующего множества.

\Exe Пусть даны два множества в СНМ, первое больше второго. Верно ли утверждение, что высота первого больше высоты второго?

\Th Пусть выполнялось $m$ операций $\texttt{MAKE\_SET}, \texttt{SAME}, \texttt{UNITE}$, из которых $n$ операций $\texttt{MAKE\_SET}$. Тогда суммарное время работы при использовании весовой эвристики равно $\O(m \log n)$.
\begin{proof}
    Вспомним, что \texttt{MAKE\_SET} работает за $\O(1)$, поэтому его суммарное время выполнения есть $\O(n)$. 
    В тоже время время основной вклад в асимтотику дают \texttt{SAME}, \texttt{UNITE}, т.к. они напрямую зависят от того на какой глубине лежат рассматриваемые элементы $x$ и $y$, ведь мы <<прыгаем>> от них к корню.
    Можно сказать, что каждая из них работает за $\Theta(h(x) + h(y))$, где $h(v)$ -- глубина, на которой лежит вершина $v$. Далее приведём оценку $h$.

    Давайте рассмотрим произвольный элемент $x$ и исследуем его положение в течение всей последовательности операций. Изначально, когда только вызвали \texttt{MAKE\_SET}, он находился на глубине $h(x) = 0$.
    Далее рассмотрим случаи объединения с некоторым множеством, которое содержит некоторый элемент $y$:
    \begin{enumerate}
        \item $s(x) > s(y) \Rightarrow$ Тогда мы подвешиваем множество с элементом $y$ к множеству с элементом $x$ и $h(x)$ \textit{не меняется}. 
        \item $s(x) < s(y) \Rightarrow$ Тогда мы подвешиваем множество с элементом $x$ к множеству с элементом $y$. Но в этом случае глубина $x$ увеличилась на единицу. Также заметим, что $s(x \cup y) > 2 s(x)$.
        Действительно, т.к. мы подвесили множество с элементом $x$ к множеству с элементом $y$, которое не менее чем первое, то мы увеличились минимум в два раза. Более формально:
        \begin{equation*}
            s(x \cup y) = s(x) + s(y) > s(x) + s(x) = 2s(x)
        \end{equation*}
        \item $s(x) = s(y) \Rightarrow$ Предположим худший случай и мы подвесили множество с $x$ к множеству с $y$ (хотя можно было и наоборот, но опять же смотрим худший случай). Тогда высота $x$ опять увеличилась на единицу и размер результирующего увеличился в два раза:
        \begin{equation*}
            s(x \cup y) = s(x) + s(y) = s(x) + s(x) = 2s(x)
        \end{equation*}
    \end{enumerate}

    Из вышеописанного следует, что высота $x$ увеличивается в худшем случае на единицу. И она увеличивается, тогда и только тогда, когда это приводит к увелечению множества, в котором лежит $x$, минимум в два раза.
    Значит, высота $x$ оценивается сверху количеством увечений его множества в два раза. А в наибольшем множестве, котором он мог находиться -- это в множестве размера $n$.
    Количество увелечений минимум в два раза до размера $n$ есть $\O(\log n)$. Значит, высота не более чем логарифмична.

    В частности из этого следует, что время работы \texttt{SAME} и \texttt{UNITE} есть $\O(\log n)$. Следовательно, в силу того, что у нас всего $m$ операций, то итоговая асимпотика есть $\O(m \log n)$.
\end{proof}

\Cons Пусть выполнялось $m$ операций $\texttt{MAKE\_SET}, \texttt{SAME}, \texttt{UNITE}$, из которых $n$ операций $\texttt{MAKE\_SET}$. Тогда амортизированная стоимость операций при использовании весовой есть $\O^*(\log n)$.
\begin{proof}

    \begin{equation*}
        \frac{\O(m \log n)}{m} = \O \left(\frac{m \log n}{m}\right) = \O^*(\log n)
    \end{equation*}
\end{proof}

\section{Ранговая эвристика\optstar}

Есть еще один способ аналогичный прошлому, его иногда используют как альтернативу. Чтобы реже приходить к ситуации бамбука и уменьшать высоту, надо подвешивать дерево с меньшей высотой к большему, потому что тогда высота почти не меняется.
Такой подход называется \textit{ранговой эвристикой}. Будем поддерживать ранг -- грубую оценку на высоту дерева. Ранг множества, в котором лежит вершина $x$, будем обозначать $r(x)$. Считается она следующим образом:
\begin{enumerate}
    \item При $\texttt{MAKE\_SET(}x\texttt{)}$ выставляем множеству ранг 0.
    \item При $\texttt{UNITE(}x, y\texttt{)}$ рассмотрим случаи:
        \begin{enumerate}
            \item $r(x) > r(y) \Rightarrow$ подвешиваем дерево, в котором лежит $y$, к дереву в котором лежит $x$. Ранг после этого равен $r(x)$ (ведь <<высота дерева>> не поменялась).
            \item $r(x) < r(y) \Rightarrow$ подвешиваем дерево, в котором лежит $x$, к дереву в котором лежит $y$. Ранг после этого равен $r(y)$ (ведь <<высота дерева>> не поменялась).
            \item $r(x) = r(y)$. Тут без разницы к кому что подвешивать, ведь ранг одинаков. Но после подвешивания ранг новоиспеченного множества будет равен $r(x) + 1$ (<<высота дерева>> увеличилась на единицу).
        \end{enumerate}
\end{enumerate}

\Exe Верно ли что если ранг одного дерева больше другого, то тоже верно для размера соответующих множеств?

\Exe Покажите, что если использовать только ранговую эвристику, то ранг есть максимальная глубина дерева.

Сама по себе эта эвристика дает уже прирост к работе. Следующие два утверждения даны для интересующегося темой читателя и принимаются без доказательства.

\ThBd Пусть выполнялось $m$ операций $\texttt{MAKE\_SET}, \texttt{SAME}, \texttt{UNITE}$, из которых $n$ операций $\texttt{MAKE\_SET}$. Тогда суммарное время работы при использовании ранговой эвристики равно $\O(m \log n)$.
% \begin{proof}

%     Вспомним, что \texttt{MAKE\_SET} работает за $\O(1)$, поэтому его суммарное время выполнения есть $\O(n)$. 
%     В тоже время время основной вклад в асимтотику дают \texttt{SAME}, \texttt{UNITE}, т.к. они напрямую зависят от того на какой глубине лежат рассматриваемые элементы $x$ и $y$, ведь мы <<прыгаем>> от них к корню.
%     Можно сказать, что каждая из них работает за $\Theta(h(x) + h(y))$, где $h(v)$ -- глубина, на которой лежит вершина $v$. Далее приведём оценку $h$.

%     Давайте рассмотрим произвольный элемент $x$ и исследуем его положение в течение всей последовательности операций. Изначально, когда только вызвали \texttt{MAKE\_SET}, он находился на глубине $h(x) = 0$.
%     Далее рассмотрим случаи объединения с некоторым множеством, которое содержит некоторый элемент $y$:
%     \begin{enumerate}
%         \item $r(x) > r(y) \Rightarrow$ Тогда мы подвешиваем множество с элементом $y$ к множеству с элементом $x \Rightarrow$ 
%     \end{enumerate}
    
% \end{proof}

\ConsBd Пусть выполнялось $m$ операций $\texttt{MAKE\_SET}, \texttt{SAME}, \texttt{UNITE}$, из которых $n$ операций $\texttt{MAKE\_SET}$. Тогда амортизированная стоимость операций $\texttt{MAKE\_SET}, \texttt{SAME}, \texttt{UNITE}$ есть $\O^*(\log n)$.
\begin{proof}

    \begin{equation*}
        \frac{\O(m \log n)}{m} = \O \left(\frac{m \log n}{m}\right) = \O^*(\log n)
    \end{equation*}

\end{proof}

То есть получается, уже сама по себе ранговая эвристика ускоряет нас с линейного до в каком-то плане логарифмического времени работы.

\section{Эвристика сжатия путей}

Также, кажется что неплохой идей будет узлы, для которых мы узнали корень (то есть все на пути при поднятии к корню), \textit{переподвешивать сразу к корню}.

Рассмотрим на примере работу данной эвристики. Пусть даны два дерева.

\begin{figure}[H]
    \centering
    \begin{tikzpicture}[
        node distance=1cm,
        state/.style={circle,draw,fill=white!20,minimum size=0.5cm,thick,font=\bfseries}
    ]

    \node[state] (b) {$b$};
    \node[state] (a) [below left=of b] {$a$};
    \node[state] (d) [below right=of b] {$d$};
    \node[state] (c) [below left=of d] {$c$};
    \node[state] (e) [below right=of d] {$e$};

    \draw[->,thick] (a) -- (b);
    \draw[->,thick] (d) -- (b);
    \draw[->,thick] (c) -- (d);
    \draw[->,thick] (e) -- (d);

    \begin{scope}[xshift=6cm]
        \node[state] (g) {$g$};
        \node[state] (f) [below left=of g] {$f$};
        \node[state] (h) [below right=of g] {$h$};
        \node[state] (i) [below right=of h] {$i$};
        
        \draw[->,thick] (f) -- (g);
        \draw[->,thick] (h) -- (g);
        \draw[->,thick] (i) -- (h);
    \end{scope}

    \end{tikzpicture}
\end{figure}

И мы хотим проверить что $e$ и $i$ лежат в одном множестве. Тогда после вызова $\texttt{SAME(}e, i\texttt{)}$ конфигурация деревьев будет следующая:

\begin{figure}[H]
    \centering
    \begin{tikzpicture}[
        node distance=1cm,
        state/.style={circle,draw,fill=white!20,minimum size=0.5cm,thick,font=\bfseries}
    ]

    \node[state] (b) {$b$};
    \node[state] (a) [below left=of b] {$a$};
    \node[state] (d) [below right=of b] {$d$};
    \node[state] (e) [right=of d] {$e$};
    \node[state] (c) [below left=of d] {$c$};

    \draw[->,thick] (a) -- (b);
    \draw[->,thick] (d) -- (b);
    \draw[->,thick] (e) -- (b);
    \draw[->,thick] (c) -- (d);

    \begin{scope}[xshift=7cm]
        \node[state] (g) {$g$};
        \node[state] (f) [below left=of g] {$f$};
        \node[state] (h) [below right=of g] {$h$};
        \node[state] (i) [right=of h] {$i$};
        
        \draw[->,thick] (f) -- (g);
        \draw[->,thick] (h) -- (g);
        \draw[->,thick] (i) -- (g);
    \end{scope}

    \end{tikzpicture}
\end{figure}

Тогда следующий вызов $\texttt{SAME(}e, i\texttt{)}$, очевидно, будет быстрее, ведь нам нужно подняться на один уровень вместо 2-х. 
Также как и ранговая, сама по себе эвристика сжатия путей тоже дает выигрыш.

\ThBd Если было выполнено $n$ операций $\texttt{MAKE\_SET}$, а также $m$ операций \texttt{SAME}. Тогда суммарное время работы при использовании ранговой эвристики равно $\O\left(n + \frac{m}{2} \log_{2 + \frac{m}{2n}} n \right)$.

\section{Объединение по рангу со сжатием путей}

Самый большой прирост даёт испольвание этих двух эвристик одновременно. Если раньше мы могли амортизированно дойти до логарифма, то с двумя эвристиками мы можем дойти до почти константы.
Давайте же рассмотрим эту <<почти константу>>.

\Def Функцией Аккермана назовём следующую функцию:
\begin{equation*}
    A(m, n) \overset{def}{=} \begin{cases}
        n + 1, m = 0 \\
        A(m - 1, 1), m > 0, n = 0\\
        A(m - 1, A(m, n - 1)), m > 0, n > 0
    \end{cases}
\end{equation*}

Функция $A$ является чрезывчайно быстро растущей функцией.

\Def Обратной функцией Аккермана положим $\alpha(n) \overset{def}{=} \min\limits_{k \in \N} \{ k: A(k, k) \geq n  \}$.

На практике $\alpha(n) \leq 4$, поэтому в рамках разумных задач, можно считать её константой (но теоретически, все же это функция).

Следующая теорема даёт нам желанный результат. Сие доказательство столь сложно, что Кормен расположил его на 10 страницах А4, поэтому, очевидно, мы принимаем его без доказательства.

\ThBd Пусть выполнялось $m$ операций $\texttt{MAKE\_SET}, \texttt{SAME}, \texttt{UNITE}$, из которых $n$ операций $\texttt{MAKE\_SET}$. 
Тогда при использовании обоих эвристик время выполнения есть $\O(m \alpha(n))$.

\ConsBd Пусть выполнялось $m$ операций $\texttt{MAKE\_SET}, \texttt{SAME}, \texttt{UNITE}$, из которых $n$ операций $\texttt{MAKE\_SET}$. Тогда амортизированная стоимость операций есть $\O^*(\alpha(n))$.
\begin{proof}
    \begin{equation*}
        \frac{\O(m \alpha(n))}{m} = \O\left(\frac{m \alpha(n)}{m}\right) = \O^*\left(\alpha(n)\right)
    \end{equation*}
\end{proof}

